%! Exponential and Logarithmic Equations and Inequalities — Unit 5, Lesson 5.6 — Guided Notes (Answer Key)
\documentclass[10pt]{article}
\usepackage{precalculus-article}
\usepackage{precalculus-key}   % pulls in -boxes; do NOT also load -boxes
\newcommand{\TallMath}[1]{$\displaystyle #1\rule[-1.4em]{0pt}{3.2em}$}

\begin{document}

\pageheader{Unit 5, Lesson 5.6}{Guided Notes --- Answer Key}
\namedateperiod

\vspace{0.08in}

\begin{objectivebox}
\textbf{LO 2.13.A:} Solve exponential and logarithmic equations and inequalities. \quad
\textbf{LO 2.13.B:} Construct the inverse function for exponential and logarithmic functions.

By the end of this lesson I will be able to:
\begin{itemize}[leftmargin=1.5em,itemsep=3pt,topsep=3pt]
  \item Solve exponential equations using same-base reasoning and by taking logarithms.
  \item Solve logarithmic equations and identify/reject extraneous solutions.
  \item Rewrite $b^x$ using the natural base identity $b^x = e^{(\ln b)\cdot x}$.
  \item Find the inverse of a transformed exponential or logarithmic function.
\end{itemize}
\end{objectivebox}

\vspace{0.06in}

\begin{vocabbox}
\termblanklong{extraneous solution}
\termblanklong{one-to-one property}
\termblanklong{natural base identity}
\termblanklong{inverse operations}
\end{vocabbox}

\vspace{0.06in}

\begin{hookbox}
\textbf{Drug half-life:} A drug has a half-life of 4 hours; initial dose is 200 mg.
When will less than 10 mg remain? We need to solve:
\[
  200\cdot\Bigl(\tfrac{1}{2}\Bigr)^{t/4} < 10
\]
What tool lets us isolate the exponent $t$?

\vspace{4pt}
My answer: \ansline{Take logarithms of both sides to bring the exponent down using the power rule.}
\end{hookbox}

\vspace{0.06in}

\begin{notesbox}{Solving Exponential Equations (EK 2.13.A.1)}

\textbf{Method 1 --- One-to-One Property (same base):}
If $b^a = b^c$ then $a = c$.

\vspace{4pt}
Example: \quad $2^x = 2^5 \;\Rightarrow\; x = $ \ans{5}

\vspace{8pt}
\textbf{Method 2 --- Take log of both sides:}

\vspace{4pt}
$3^x = 15$

\vspace{4pt}
Step 1: Take $\log$ of both sides: \quad $\log(3^x) = \log 15$

\vspace{4pt}
Step 2: Power rule: \quad $x \cdot \log 3 = \log 15$

\vspace{4pt}
Step 3: Divide: \quad $x = \dfrac{\log 15}{\log 3} \approx $ \ans{2.465}

\vspace{10pt}
\textbf{Your turn:} Solve $5^x = 100$. Show all steps.

\vspace{6pt}
Step 1: \ansline{Take log of both sides: $\log(5^x) = \log 100$}

\vspace{4pt}
Step 2: \ansline{Power rule: $x \cdot \log 5 = \log 100 = 2$}

\vspace{4pt}
Step 3: $x = $ \ans{$2/\log 5 \approx 2.861$}

\end{notesbox}

\vspace{0.06in}

\begin{notesbox}{Solving Logarithmic Equations + Checking Extraneous Solutions (EK 2.13.A.1--2)}

\textbf{Convert to exponential form:}

\vspace{4pt}
$\log_2 x = 4 \;\Rightarrow\; x = 2^4 = $ \ans{16} \quad Check: $\log_2($ \ans{16} $) = 4$? \ans{Yes}

\vspace{8pt}
$\log_3(x+2) = 2 \;\Rightarrow\; x+2 = 3^2 = $ \ans{9} $\;\Rightarrow\; x = $ \ans{7}

\vspace{10pt}
\textbf{Combined logs --- check for extraneous solutions:}

\vspace{4pt}
Solve: $\log_5 x + \log_5(x-4) = 1$

\vspace{4pt}
Step 1 (combine): $\log_5\bigl(x(x-4)\bigr) = 1$

\vspace{4pt}
Step 2 (exponential form): $x(x-4) = $ \ans{5}

\vspace{4pt}
Step 3 (expand and solve): $x^2 - 4x - $ \ans{5} $= 0 \;\Rightarrow\; (x - $ \ans{5} $)(x + $ \ans{1} $) = 0$

\vspace{4pt}
Possible solutions: $x = $ \ans{5} or $x = $ \ans{$-1$}

\vspace{4pt}
Check domain (argument must be $> 0$): \ansline{$x=-1$ gives $\log_5(-1)$, undefined. Extraneous. Reject $x=-1$.}

\vspace{4pt}
Final answer: $x = $ \ans{5}

\end{notesbox}

\vspace{0.06in}

\begin{notesbox}{Natural Base Identity (EK 2.13.A.3)}

\[
  b^x = e^{(\ln b)\cdot x}
\]

\vspace{4pt}
This means every exponential function can be written using base $e$.

\vspace{6pt}
Fill in:
\begin{itemize}[leftmargin=1.5em,itemsep=6pt,topsep=4pt]
  \item $2^x = e^{({\color{keyred}\mathbf{\ln 2}})\cdot x} \approx e^{0.693x}$
  \item $10^x = e^{({\color{keyred}\mathbf{\ln 10}})\cdot x} \approx e^{2.303x}$
  \item $\bigl(\tfrac{1}{2}\bigr)^x = e^{({\color{keyred}\mathbf{\ln(1/2)}})\cdot x} \approx e^{-0.693x}$
\end{itemize}

\end{notesbox}

\vspace{0.06in}

\begin{notesbox}{Inverse of a Transformed Exponential (EK 2.13.B.1)}

Find the inverse of $f(x) = 3 \cdot 2^{x-1} + 4$.

\vspace{4pt}
Step 1: Replace $f(x)$ with $y$: \quad $y = 3 \cdot 2^{x-1} + 4$

\vspace{4pt}
Step 2: Swap $x$ and $y$: \quad $x = 3 \cdot 2^{y-1} + 4$

\vspace{4pt}
Step 3: Subtract 4: \quad $x - 4 = $ \ans{$3 \cdot 2^{y-1}$}

\vspace{4pt}
Step 4: Divide by 3: \quad $\dfrac{x-4}{3} = $ \ans{$2^{y-1}$}

\vspace{4pt}
Step 5: Apply $\log_2$: \quad $\log_2\!\dfrac{x-4}{3} = $ \ans{$y - 1$}

\vspace{4pt}
Step 6: Add 1: \quad $f^{-1}(x) = $ \ans{$\log_2\!\dfrac{x-4}{3} + 1$}

\end{notesbox}

\vspace{0.06in}

\begin{notesbox}{Inverse of a Transformed Logarithm (EK 2.13.B.2)}

Find the inverse of $f(x) = 2\log_3(x+1) - 5$.

\vspace{4pt}
Step 1: $y = 2\log_3(x+1) - 5$; swap $x \leftrightarrow y$: $\;x = 2\log_3(y+1) - 5$

\vspace{4pt}
Step 2: Add 5: \quad $x + 5 = $ \ans{$2\log_3(y+1)$}

\vspace{4pt}
Step 3: Divide by 2: \quad $\dfrac{x+5}{2} = $ \ans{$\log_3(y+1)$}

\vspace{4pt}
Step 4: Exponentiate (base 3): \quad $3^{(x+5)/2} = $ \ans{$y + 1$}

\vspace{4pt}
Step 5: Subtract 1: \quad $f^{-1}(x) = $ \ans{$3^{(x+5)/2} - 1$}

\end{notesbox}

\end{document}
